%% ============================================================
%%  生成AI利用申告書 — テンプレート
%%  ハッカソン1 提出課題用 (塩澤 / 2026)
%%  コンパイル: docker run --rm -v "$(pwd):/workspace" -w /workspace \
%%              ghcr.io/paperist/texlive-ja:debian latexmk main.tex
%% ============================================================
\documentclass[a4paper,11pt,dvipdfmx]{jsarticle}

\usepackage[top=25mm,bottom=25mm,left=25mm,right=25mm]{geometry}
\usepackage{amssymb}
\usepackage{listings}
\usepackage{plistings}  % 日本語コメント混在時の文字幅補正
\usepackage[dvipdfmx,hidelinks]{hyperref}
\usepackage{pxjahyper}

\lstset{
  basicstyle=\ttfamily\small,
  frame=single,
  breaklines=true,
  xleftmargin=1em,
  aboveskip=6pt, belowskip=6pt,
  columns=flexible,
}

\newcommand{\cboxon}{$\boxtimes$}
\newcommand{\cboxoff}{$\square$}

\begin{document}

{\LARGE\bfseries 生成AI利用申告書}
\vspace{1em}

\begin{itemize}
  \item 氏名：佐藤涼菜
  \item 学籍番号：25G1059
  \item プログラムファイル名： HCK02\_01.ino
  \item 使用したAIツール（例：ChatGPT など）： Gemini pro, claude
\end{itemize}

\section*{1. 何に使ったか}
今回の課題で、生成AIをどの部分に使いましたか。


提出課題Ⅰの内容理解とコードの案だし．

\section*{2. AIへの指示（プロンプト）}
AIにどのような質問・指示を出しましたか。重要なものを1〜2個書いてください。

・サンプリング周期を固定化するためのコード案を出すこと\par
・出力される波形が安定するためにはどのような要素を改善，削除すべきか

\section*{3. AIの出力の使い方}
AIの出力をどのように使いましたか。
\begin{itemize}
  \item \cboxon 一部修正して使った
  \item \cboxon 参考にしただけ
\end{itemize}
$\rightarrow$ 上で選んだ理由を書いてください：
サンプリング周期を安定させるためのコード案においては，提案されたコードの内容を理解した上で,実機での動作検証や他の生成AIによるダブルチェックを行い，必要に応じて修正を加えたため．
また，課題の内容が正しく認識できているか確認するために生成AIを利用したため，内容の齟齬がみられなければそのまま課題に
取り組んだので参考にしただけを選択した．

\section*{4. 正しいと確認した方法}
AIの出力が正しいと、どのように確認しましたか。


実際に実行し，動作が想定の通りに実行されたかどうかを確認し，コードを読んで意味を理解するまで調べた．また，関数や変数の仕組みについては，
Arduino日本語リファレンスで確認をしハルシネーションが発生してないか確認した．更に，
別の生成AIにも同様のプロンプトを送り，回答結果を検証することで，信憑性を高めた．


\section*{5. 問題や間違い}
AIの出力に問題や間違いはありましたか。
\begin{itemize}
  \item \cboxon あった
  \item \cboxoff なかった（※本当に？と思って確認した内容を書くこと）
\end{itemize}
確認した内容や気づいたこと：
同じプロンプトで質問をしても，生成AIの種類によって出力結果は異なっていた．また，何度も同様の質問を行うことにより，
出力内容の精度が向上したり，逆に間違った解釈をされてしまったりと一概に出力結果を鵜呑みにしてはいけないことがわかった．
新しいチャットに同様のプロンプトを投げることや，別の生成AIを用いて検証を行うことが
ハルシネーションを鵜呑みにしない方法であると考えた．

\section*{6. 自分で工夫した点}
例題２では出力結果に最大振幅や測定値など，漢字でキャプションを書いていたが，これらはより正確な波形を出力するために
は不要であったため削除した．また，読み取ったアナログ値を直接出力することにより，余計な計算が省かれprocessingが描画する正弦波に
近しい波形を出力することができた．

\section*{7. 参考文献}

  ・Arduino日本語リファレンス.  
  \url{https://www.musashinodenpa.com/arduino/ref/}\par

  ・Arduino公式ドキュメント
  \url{https://docs.arduino.cc/}\par



 
\end{document}
